% Chapter 2 -- Literature Review
% Status: basic draft built from Papers/digested/*.md. The comparison table of
% published numbers (Table 2.1) must be checked cell by cell against the digests
% before submission; cells marked \pending were not verified.
\chapter{Literature Review}
\label{ch:literature}

\section{Collaborative Filtering with Implicit Feedback}
\label{sec:lit-cf}

Matrix factorisation (MF) represents user $u$ and item $i$ by vectors $\pu, \qi \in
\mathbb{R}^d$ and predicts $\hat{s}_{ui} = \pu^\top \qi$ plus bias terms. With implicit feedback
(clicks, views, ratings treated as presence) the training signal is pairwise or pointwise: BPR
\cite{rendle2009bpr} maximises $\log\sigma(\hat{s}_{ui} - \hat{s}_{uj})$ for an observed item
$i$ and a sampled unobserved item $j$; the pointwise alternative treats observed pairs as
positives and sampled unobserved pairs as negatives under a binary cross-entropy (BCE) loss, the
choice of neural collaborative filtering (NCF) \cite{he2017ncf} and of PFedRec
\cite{zhang2023pfedrec}. Coordinate-descent and graph-based variants
\cite{bayer2016implicit,he2020lightgcn} change the optimiser or the encoder but keep the
inner-product scorer.

The evaluation protocol of NCF is the de facto standard for these models: the latest
interaction of each user is held out, 99 unobserved items are sampled as candidates, and \hr{}
and \ndcg{} measure where the held-out item ranks among the 100 candidates. Krichene and Rendle
\cite{krichene2020sampled} show that sampled metrics are inconsistent with their exhaustive
counterparts and are sensitive to the sampling scheme, which is one reason why numbers from
different code bases are not directly comparable; this thesis therefore scores every method
with the same candidate lists per user.

\section{Federated Learning}
\label{sec:lit-fl}

FedAvg \cite{mcmahan2017fedavg} trains a global model by repeating: sample a fraction $C$ of
the $N$ clients, send them the global parameters, let each run $K$ local epochs, and average
the returned parameters weighted by local data size. Heterogeneous data make the local updates
drift apart; FedProx \cite{li2020fedprox} adds a proximal term $\frac{\mu}{2}\|\theta -
\vect{w}\|^2$ to the local objective, SCAFFOLD \cite{karimireddy2020scaffold} corrects the drift
with control variates, and adaptive server optimisers \cite{reddi2021adaptive} replace the
plain average by Adam-like updates. Kairouz et al.\ \cite{kairouz2021advances} distinguish the
\emph{cross-silo} setting (few, reliable, stateful clients) from the \emph{cross-device} setting
(many unreliable clients, each participating rarely, stateless or with limited state). The
protocol of this thesis is cross-device: $N=\NClients$ clients, $C=\ParticipationC$, and
client state that persists between participations.

\section{Personalised Federated Learning}
\label{sec:lit-pfl}

A single global model serves heterogeneous clients poorly, which motivates personalised FL.
Three families are relevant here.

\paragraph{Parameter splitting.} FedPer \cite{arivazhagan2019fedper} shares the base layers and
keeps the top layers private; FedRep \cite{collins2021fedrep} shares a representation and keeps a
low-dimensional head per client, updating the head for many more steps than the body. Federated
Reconstruction \cite{singhal2021fedrecon} goes further for stateless clients: the private part
(for MF, the user vector) is reconstructed from scratch at every participation and at inference
with the global part frozen, which the authors deployed in Gboard.

\paragraph{Interpolation.} APFL \cite{deng2020apfl} forms the personal model as $\alpha_u
\theta_u + (1-\alpha_u)\vect{w}$ and learns $\alpha_u$ by gradient descent on the local loss;
its analysis gives a closed-form optimal $\alpha_u$ that grows with the divergence between the
local and global optima. Mansour et al.\ \cite{mansour2020three} compare clustering,
interpolation and fine-tuning as three routes to personalisation.

\paragraph{Test-time personalisation.} FedTHE \cite{jiang2023fedthe} keeps a global head and a
local head on a shared extractor and chooses a per-sample ensemble weight at test time by
entropy minimisation; its gains appear mainly under distribution shift at test time and are
small in-distribution. This is the closest prior work to the blending experiment of
Section~\ref{sec:res-placement}; following the advisor's guidance, the results of this thesis are
described as \emph{consistent with} the findings of APFL and FedTHE rather than as confirming
them.

\section{Federated Recommendation}
\label{sec:lit-fedrec}

Early federated CF kept the user vector on the device and aggregated the item vectors
\cite{ammaduddin2019fcf,chai2021fedmf}, the placement called \emph{Personalized} in this
thesis. The methods below personalise more.

\paragraph{PFedRec} \cite{zhang2023pfedrec} removes the user vector altogether: each client
holds a private scoring function (a one-layer head $\wu$) on top of a shared item table
$\Qtab$, trained by alternating updates of the head and of the item rows. At inference the
official code scores each user with the client's own \emph{post-tuned} copy of the item table
from the last local epoch, not with the server's table. The paper reports
\PublishedPfedrecNdcg{} \ndcg{} and \PublishedPfedrecHr{} \hr{} on ML-1M with full
participation ($C=1$), 100 rounds, one local epoch and a validation item (the second-latest
interaction). Its released code excludes every interacted item, including the held-out test
item, from the training negatives.

\paragraph{FedRAP} \cite{li2024fedrap} adds an additive private item table $\vect{D}_u$ to a
sparse global one, with curriculum regularisers; its released code selects the best round on the
test metric and its published ML-1M numbers are therefore not used as a target here.
\paragraph{GPFedRec} \cite{zhang2024gpfedrec} aggregates item tables per client along a
server-side user-similarity graph; its own PFedRec re-run on ML-1M (\pending{73.03/44.49, verify
in the digest}) is above its proposed method. \paragraph{CoFedRec} \cite{he2024cofedrec} clusters
clients by item category and adds supervised contrastive learning; \paragraph{FedCA}
\cite{zhang2024fedca} combines similarity and complementarity in aggregation; \paragraph{LightFR}
\cite{zhang2022lightfr} learns binary codes for communication efficiency; \paragraph{P2FedRec}
\cite{hu2025p2fedrec} adds secret sharing, homomorphic encryption and local differential privacy
on top of relationship-aware aggregation. \paragraph{FedDAE} \pending{digest of FedDAE (AAAI
2025) once the PDF is available; add to Table~\ref{tab:lit-numbers} and to the comparison in
Section~\ref{sec:lit-positioning}}. Surveys of the area
\cite{zhang2025survey,yin2025devicers} organise these methods by which part is personalised
and by the aggregation rule, which is the organisation adopted in Chapter~\ref{ch:method}.

\begin{table}[htbp]
  \caption{Published ML-1M numbers of related methods with the protocol details that affect
  comparability. ``Own table'' = the client scores with its own post-tuned item table;
  ``test excluded'' = the held-out test item is removed from the training negatives. Numbers are
  as reported in the papers or in their re-runs; every cell must be re-checked against
  \code{Papers/digested/} before submission.}
  \label{tab:lit-numbers}
  \centering\small
  \begin{tabular}{lrrllll}
    \toprule
    Method & \hr{} & \ndcg{} & $C$ & Own table & Validation item & Test excluded \\
    \midrule
    NCF / NeuMF (centralised) \cite{he2017ncf} & 0.730 & 0.447 & -- & -- & yes & \pending{?} \\
    FedRecon (PFedRec re-run) \cite{singhal2021fedrecon,zhang2023pfedrec} & 0.633 & 0.366 & 1.0 & no & yes & yes \\
    PFedRec \cite{zhang2023pfedrec} & \PublishedPfedrecHr & \PublishedPfedrecNdcg & 1.0 & yes & yes & yes \\
    GPFedRec \cite{zhang2024gpfedrec} & 0.722 & 0.436 & 1.0 & yes & yes & \pending{?} \\
    CoFedRec \cite{he2024cofedrec} & 0.778 & 0.488 & \pending{?} & \pending{?} & \pending{?} & \pending{?} \\
    FedRAP \cite{li2024fedrap} & 0.932 & 0.719 & 1.0 & yes & no (best round on test) & yes \\
    LightFR \cite{zhang2022lightfr} & 0.501 & -- & \pending{?} & no & \pending{?} & \pending{?} \\
    \midrule
    This thesis, PFedRec (Section~\ref{sec:res-pfedrec}) & \HrPfedrec & \NdcgPfedrec & \ParticipationC & no & no & no (clean) \\
    \bottomrule
  \end{tabular}
\end{table}

\section{Privacy in Federated Recommendation}
\label{sec:lit-privacy}

Federated training does not by itself protect the user. Gradients and parameter updates leak
training data \cite{zhu2019dlg,melis2019exploiting}; in MF the update of the item table is
non-zero only for the items a user interacted with, so an honest-but-curious server learns the
interaction set from the sparsity pattern alone \cite{chai2021fedmf}. Uploading the user vector
is worse: with the item table known to the server, $\pu^\top\Qtab^\top$ ranks every item for the
user and recovers the preferences directly, which is why Federated Reconstruction never
transmits nor persists the user vector server-side \cite{singhal2021fedrecon}. Mitigations are
secure aggregation \cite{bonawitz2017secagg}, which hides individual updates behind the sum,
local differential privacy and cryptographic aggregation \cite{hu2025p2fedrec}, and binary codes
that resist inversion \cite{zhang2022lightfr}. GPFedRec and FedRAP apply Laplace or Gaussian
noise but do not report the resulting $\varepsilon$ \cite{zhang2024gpfedrec,li2024fedrap}. The
threat model of Section~\ref{sec:method-threat} and the attack of Section~\ref{sec:res-privacy}
follow this line: the server is honest-but-curious, and the attack is exactly the
$\pu^\top\Qtab^\top$ ranking that the adaptive method's upload of $\pu$ makes possible.

\section{Evaluation Pitfalls in Federated Recommendation}
\label{sec:lit-pitfalls}

Four protocol choices change the reported numbers by amounts comparable to the differences
between methods.
\begin{enumerate}[leftmargin=2em]
  \item \textbf{Participation fraction.} Most FedRec papers use $C=1$; cross-device deployments
        use $C \ll 1$, under which the private part of a client is stale between participations
        \cite{kairouz2021advances}.
  \item \textbf{Inference with a client-side item table.} PFedRec and FedRAP score with a
        post-tuned or additive private item table; methods that score with the shared table are
        at a disadvantage that has nothing to do with the placement of personalisation.
  \item \textbf{Validation item.} Selecting the best round on the test item inflates the result;
        PFedRec and FedRAP hold out a second item for validation, many re-implementations do not.
  \item \textbf{Test item in the negative pool.} Excluding every interacted item, including the
        held-out test item, from the training negatives is common in released code. With a shared
        table the leak is diluted across users; with any per-user fine-tuning of the item table it
        is not (Section~\ref{sec:res-pfedrec}).
\end{enumerate}
To these the cross-device setting adds a fifth, specific to stateful clients: the evaluation
must load each client's trained local state, and a scoring pass that silently uses a cold
(initial) state produces a plausible-looking but meaningless number
(Section~\ref{sec:setup-validity}).

\section{Positioning of this Thesis}
\label{sec:lit-positioning}

The methods above each propose a placement and report it under its own protocol. This thesis
holds the protocol fixed (cross-device, $C=\ParticipationC$, shared candidate lists, clean
negatives, no validation item declared as such) and varies the placement: all-global, private
user vector, private head, and blend with population knowledge. It measures the protocol
effects that the literature leaves implicit (post-tuning with and without the test-item leak;
staleness at low participation) and reports the privacy cost of each placement in the same
threat model. The adaptive method of Chapter~\ref{ch:method} is the thesis's own design and is
reported as a negative result in the sense of the FedTHE observation that blending helps mainly
under test-time shift.
